% Hafta 3 — HKDF: bir ana sırdan birçok anahtar
% Derleme: py -3.12 tools/latex/derle.py docs/week-3/assets/h03-03-hkdf.tex
\documentclass[tikz,border=8pt]{standalone}
\usepackage{cen429-cizim}
\begin{document}
\begin{tikzpicture}[node distance=8mm and 16mm]
  % Başlık
  \node[baslik] (b) at (0,0) {HKDF: bir ana sırdan birçok anahtar};
  \node[altbaslik] at (b.south west) {Extract (topla) $\rightarrow$ Expand (amaca göre böl)};

  % Ana zincir: sır -> extract -> PRK
  \node[kutu=24mm] (sir) at (0.5,-2.4) {\kb{Ana sır}};
  \node[koyukutu=28mm, right=of sir] (extract) {HKDF-Extract\\+ tuz};
  \node[kutu=18mm, fill=acik, right=of extract] (prk) {\kb{PRK}};
  \draw[ok, draw=soluk] (sir.east) -- (extract.west);
  \draw[ok, draw=soluk] (extract.east) -- (prk.west);

  % PRK'dan üç türetilmiş anahtar
  \node[iyikutu=34mm] (k1) at ($(prk)+(5.8,2.4)$) {\kb{Şifreleme anahtarı}\\\kod{info='sifreleme'}};
  \node[iyikutu=34mm] (k2) at ($(prk)+(5.8,0)$) {\kb{MAC anahtarı}\\\kod{info='mac'}};
  \node[iyikutu=34mm] (k3) at ($(prk)+(5.8,-2.4)$) {\kb{Oturum-42 anahtarı}\\\kod{info='oturum-42'}};
  \draw[iyiok] (prk.east) -- (k1.west);
  \draw[iyiok] (prk.east) -- (k2.west);
  \draw[iyiok] (prk.east) -- (k3.west);

  % Dip şeridi
  \node[dip, text width=160mm, below=10mm of k3.south -| extract, anchor=north] {%
    Farklı \kod{info} $\rightarrow$ farklı anahtar. Aynı sırdan üretilseler de biri sızarsa diğerleri güvende kalır.};
\end{tikzpicture}
\end{document}
